\documentclass[12pt,a4paper]{report}
\usepackage[utf8]{inputenc}
\usepackage[T1]{fontenc}
\usepackage[spanish,es-nodecimaldot]{babel}
\usepackage{lmodern}
\usepackage[a4paper, left=3.5cm, right=2.5cm, top=2.5cm, bottom=2.5cm]{geometry}
\usepackage{graphicx}
\usepackage{setspace}
\usepackage{booktabs}
\usepackage{longtable}
\usepackage{array}
\usepackage{enumitem}
\usepackage{csquotes}
\usepackage[hidelinks]{hyperref}
\usepackage{float}

\onehalfspacing
\setlength{\parindent}{1.25cm}
\setlength{\parskip}{0pt}
\setlist[enumerate]{leftmargin=*}
\setlist[itemize]{leftmargin=*}

\addto\captionsspanish{\renewcommand{\tablename}{Tabla}}
\addto\captionsspanish{\renewcommand{\listtablename}{Índice de tablas}}

\begin{document}

\begin{titlepage}
\begin{center}
\vspace*{1cm}
{\large UNIVERSIDAD PRIVADA DOMINGO SAVIO}\\[0.4cm]
{\large Facultad de Ingeniería}\\[0.3cm]
{\large Ingeniería de Sistemas}\\[2cm]

{\bfseries\Large DISEÑO Y DESARROLLO DE UN SISTEMA DE GESTIÓN Y TRAZABILIDAD DE SERVICIOS DE SOPORTE TÉCNICO CON ASISTENCIA MEDIANTE INTELIGENCIA ARTIFICIAL PARA LOS PRESTADORES DE SERVICIOS DE LA CIUDAD DE TARIJA DURANTE LAS GESTIONES 2026--2027}\\[2cm]

{\large PERFIL DE PROYECTO DE GRADO}\\[2cm]

\begin{flushright}
Autor: Jesus Rafael Aviles Poita\\[0.5cm]
Tutor: Fridel Omar Baldiviezo Osina
\end{flushright}

\vfill
{\large Tarija -- Bolivia}\\
{\large 2026}
\end{center}
\end{titlepage}
\chapter*{RESUMEN}
\addcontentsline{toc}{chapter}{RESUMEN}

El presente perfil de proyecto de grado tiene como propósito fundamentar y estructurar el diseño y desarrollo de un sistema informático de gestión y trazabilidad de servicios de soporte técnico con asistencia mediante inteligencia artificial, orientado a prestadores de servicios informáticos en la ciudad de Tarija durante las gestiones 2026--2027. La investigación parte de la necesidad de analizar cómo los prestadores de servicios de soporte técnico de la ciudad de Tarija gestionan la recepción de equipos, el registro de diagnósticos, la confirmación de la intervención técnica y el seguimiento de las reparaciones. Mediante el trabajo de campo se identificarán las modalidades de registro utilizadas y las posibles dificultades de comunicación y consulta de información por parte de los clientes, con el propósito de fundamentar los requerimientos del sistema propuesto. La metodología adopta un enfoque mixto (cualitativo y cuantitativo), de nivel descriptivo-propositivo y diseño no experimental transversal, aplicando un muestreo no probabilístico por conveniencia sobre dos unidades de análisis: 10 a 15 técnicos o propietarios de talleres y 50 a 100 usuarios. La solución propuesta se adscribe al área de la ingeniería de software bajo la arquitectura Modelo-Vista-Controlador (MVC) con el framework PHP Laravel, MySQL y una interfaz responsiva con Bootstrap 5, donde el diseño contempla un módulo de trazabilidad mediante estados secuenciales, evidencias fotográficas y restricciones de borrado en base de datos. Asimismo, se propone integrar un asistente conversacional basado en inteligencia artificial estructurado sobre una API REST de un proveedor externo, tomando Google Gemini como alternativa tecnológica de referencia, cuya selección definitiva se realizará considerando su disponibilidad, continuidad del servicio, latencia de procesamiento y capacidad para procesar consultas en español. El backend de Laravel validará las solicitudes, consultará únicamente los datos autorizados y proporcionará al asistente el contexto necesario para responder preguntas frecuentes y consultas de estado mediante código verificado como característica prevista en la solución. La propuesta se fundamenta en marcos de gobernanza como ITIL y estándares ISO/IEC 25010 e ISO/IEC 27001, y contempla el desarrollo de un prototipo funcional (MVP) evaluado según criterios de calidad de software.
\chapter*{ABSTRACT}
\addcontentsline{toc}{chapter}{ABSTRACT}

This degree project proposal aims to ground and structure the design of a technical support management and traceability computer system powered by artificial intelligence, tailored for computer service providers in the city of Tarija during the 2026 term. The research stems from the need to analyze how technical support service providers in the city of Tarija manage equipment intake, diagnostic records, budget communication, and repair tracking. Through field research, the study will identify the recording methods currently used and potential communication and information access difficulties experienced by clients, in order to ground the requirements of the proposed system. The methodology adopts a mixed-methods approach (qualitative and quantitative), with a descriptive-propositive level and a non-experimental cross-sectional design, utilizing a non-probabilistic convenience sampling strategy across two analysis units: 10 to 15 workshop technicians or owners and 50 to 100 computer service clients. The proposed software engineering solution adopts a Model-View-Controller (MVC) architecture utilizing the PHP Laravel framework, MySQL, and a responsive Bootstrap 5 interface, where the design contemplates a traceability module featuring sequential service statuses, photographic evidence, and soft-delete database constraints. Furthermore, it proposes integrating a large language model (LLM) conversational assistant through an external provider REST API, taking Google Gemini as a reference technological alternative, with the final model selection subject to service availability, continuity, usage costs, and Spanish processing capability. The Laravel backend will validate requests, query strictly authorized data, and provide the assistant with the necessary context to handle frequent inquiries and status lookups as a designed architectural capability. The proposal is grounded in IT governance frameworks such as ITIL and international standards including ISO/IEC 25010 and ISO/IEC 27001.

\pagenumbering{roman}
\tableofcontents
\clearpage
\pagenumbering{arabic}

\chapter{Introducción}
\section{Antecedentes}

El estudio del estado del arte constituye un requisito metodológico fundamental para contextualizar la problemática de la gestión y trazabilidad de los servicios de soporte técnico. A continuación, se presentan las investigaciones verificadas en el ámbito internacional, organizadas de forma cronológica ascendente bajo un análisis estructurado en tres momentos: contextualización temática, síntesis parafraseada bajo norma APA 7.\textsuperscript{a} con precisión metodológica y tecnológica, y valoración del aporte específico al presente proyecto de grado.

\subsection{Antecedentes internacionales}

\subsubsection{Investigación 1: \enquote{Helpdesk como alternativa de soporte y gestión tecnológica para una empresa tabasqueña} (Ortiz González, Gómez Lemus y Zurita Cruz, 2021)}

En el ámbito internacional latinoamericano, la estandarización de procesos de atención técnica mediante marcos de gobernanza de tecnologías de información y plataformas de mesa de ayuda representa una estrategia consolidada para superar la desorganización operativa y la falta de canales estructurados de comunicación.

Ortiz González, Gómez Lemus y Zurita Cruz (2021) desarrollaron una investigación orientada a diseñar e implementar una mesa de ayuda (\textit{Helpdesk}) basada en las mejores prácticas de la biblioteca ITIL (\textit{Information Technology Infrastructure Library}) para la empresa de servicios informáticos Techfusione, en Tabasco, México. La metodología articuló las fases de estrategia, diseño, transición y operación de servicios, evaluando la calidad del software mediante métricas de funcionalidad, usabilidad y eficiencia sobre la plataforma ManageEngine ServiceDesk Plus. El estudio reportó la factibilidad de estructurar catálogos de servicios, estandarizar tiempos de respuesta por prioridad y formalizar los flujos de atención desde la recepción hasta la posventa.

\paragraph{Aporte al proyecto.}
El valor sustantivo de este antecedente radica en la aplicación del marco ITIL para catalogar los servicios de mantenimiento e identificar los flujos formales de atención técnica. La presente investigación toma como referencia los criterios de clasificación de incidentes y requerimientos para diseñar la arquitectura del software en Tarija, diferenciándose al incorporar un asistente conversacional con inteligencia artificial enfocado en la interacción directa y automatizada con el cliente.

\subsubsection{Investigación 2: \enquote{Desarrollo de una plataforma web para la gestión de órdenes de servicio de la empresa Computech} (Mosquera Perea, Mazo Agudelo y Martínez Herrera, 2025)}

La transición desde registros informales en papel o mensajería desestructurada hacia plataformas de software web centralizadas constituye una tendencia indispensable para garantizar la persistencia del historial de mantenimiento y la trazabilidad de los procesos de reparación.

Mosquera Perea, Mazo Agudelo y Martínez Herrera (2025), en la Universidad Católica de Pereira (Colombia), desarrollaron una aplicación web orientada a la gestión y seguimiento de órdenes de servicio para la empresa de soporte Computech. Construida bajo la arquitectura Modelo-Vista-Controlador (MVC), utilizando el framework PHP Laravel, MySQL, Bootstrap 5 y la metodología ágil Scrum, la plataforma implementó un esquema de seguridad basado en roles (administrador y técnico) y restricciones de integridad que prohíben la eliminación física de registros. Según la evaluación del sistema (p. 56 del documento original), la búsqueda de información e historial de servicios se redujo en un 60\%, optimizando la eficiencia operativa del taller.

\paragraph{Aporte al proyecto.}
Este trabajo aporta de manera directa el modelo entidad-relación para la gestión de órdenes de trabajo (vinculación entre cliente, dispositivo, diagnóstico y estado de reparación) y la regla de negocio de no eliminación para garantizar una trazabilidad inalterable. El presente proyecto adopta esta estructura de persistencia relacional, pero la extiende al integrar un módulo de asistencia con inteligencia artificial para la consulta remota por parte del usuario final.

\subsection{Síntesis conclusiva del estado del arte}

La revisión de los antecedentes internacionales evidencia que la estandarización de los servicios de soporte técnico mediante software web y marcos de trabajo como ITIL puede reducir los tiempos de búsqueda de información y formalizar la relación con el cliente. No obstante, se identifica una oportunidad de innovación al combinar estas arquitecturas relacionales de trazabilidad con asistentes virtuales conversacionales impulsados por inteligencia artificial, constituyendo el núcleo del diseño propuesto para la ciudad de Tarija.

\section{Delimitación}

\subsection{Límite Sustantivo}

El proyecto se adscribe al área de la Ingeniería de Software, abarcando el análisis de requerimientos, el modelado de los flujos de trabajo de las órdenes de servicio, el diseño de la arquitectura del software y el desarrollo del prototipo funcional (MVP) de la solución web con asistencia mediante inteligencia artificial. No incluye la implementación comercial física en producción en los talleres locales ni la integración con sistemas externos de facturación o inventarios corporativos.

\subsection{Límite Temporal}

El proyecto comprende las gestiones 2026 y 2027, abarcando la fase de perfil, análisis y diseño arquitectónico durante la gestión 2026 (Taller de Grado I), y la fase de desarrollo del prototipo, verificación de requerimientos y evaluación de calidad del software durante la gestión 2027 (Taller de Grado II).

\subsection{Límite Geográfico}

El ámbito de estudio se circunscribe a los prestadores de servicios de soporte técnico de equipos informáticos ubicados en el radio urbano de la ciudad de Tarija.
\section{Planteamiento del problema}

En la prestación de servicios de soporte técnico y mantenimiento de equipos informáticos, la gestión documental de la recepción de dispositivos, el registro detallado del diagnóstico técnico y la comunicación continua sobre el avance de las órdenes de trabajo constituyen actividades operativas críticas para garantizar la transparencia y la calidad en la atención al cliente.

En el contexto de los prestadores de servicios de soporte técnico de la ciudad de Tarija, la ausencia de herramientas informáticas especializadas conduce a que las actividades de recepción y seguimiento se realicen habitualmente de manera empírica e informal. En la práctica cotidiana, los datos de los equipos recibidos y los diagnósticos iniciales suelen registrarse en boletas físicas de papel o en conversaciones dispersas a través de aplicaciones de mensajería instantánea. Esta dispersión documental dificulta la reconstrucción del historial de mantenimiento, complica la verificación de los componentes sustituidos y limita la capacidad de auditar las acciones realizadas entre el técnico y el usuario.

Asimismo, los clientes experimentan barreras para consultar de forma oportuna el estado actualizado de su reparación, dependiendo de llamadas telefónicas o visitas presenciales que interrumpen la labor técnica del personal en taller. Esta dinámica genera incertidumbre en los usuarios respecto a los tiempos de entrega y la fidelidad de los diagnósticos, planteando como hipótesis de trabajo la existencia de una desconfianza latente que requiere ser investigada empíricamente en el campo mediante entrevistas y encuestas estructuradas.
\section{Formulación del problema}

¿Cómo diseñar un sistema de gestión y trazabilidad de servicios de soporte técnico con asistencia mediante inteligencia artificial, de acuerdo con los requerimientos identificados en los prestadores de servicios de la ciudad de Tarija durante la gestión 2026?
\section{Sistematización del problema y abordaje de la solución}

Para articular el problema investigado con los componentes de la propuesta de ingeniería de software, se presenta la siguiente estructura de correspondencia funcional:

\begin{table}[H]
\centering
\small
\renewcommand{\arraystretch}{1.3}
\begin{tabular}{|p{4cm}|p{4.5cm}|p{4.5cm}|}
\hline
\textbf{Situaciones preliminares bajo investigación (Diagnóstico)} & \textbf{Problema Central y Componente de Solución Propuesta} & \textbf{Consecuencias e Impactos Esperados} \\
\hline
C1: Posible uso predominante de registros manuales o canales de mensajería dispersos para documentar la recepción y los avances técnicos. & A1: Módulo de recepción digitalizada con registro de evidencias físicas (fotografías, observaciones iniciales) y emisión de comprobante único. & F1: Facilitar la verificación transparente del estado de entrega del equipo y brindar respaldo documental claro al cliente y al personal técnico. \\
\hline
C2: Dificultad del cliente para consultar el avance actualizado del mantenimiento sin requerir la interrupción del trabajo del técnico. & A2: Módulo de trazabilidad basado en estados secuenciales del servicio (recepción, diagnóstico, reparación, pruebas, entrega), marcas temporales (timestamps), registro auditable de cambios y restricciones de eliminación en base de datos. & F2: Permitir la consulta organizada de las etapas de atención con registro auditable de cambios, reduciendo la incertidumbre del usuario y las consultas repetitivas. \\
\hline
C3: Demanda de atención informativa repetitiva sobre el funcionamiento del servicio, presupuestos o interpretación de estados de órdenes de trabajo. & A3: Asistente conversacional con Inteligencia Artificial enfocado en responder preguntas frecuentes y consultar estados mediante código de orden verificado. & F3: Agilizar la atención informativa al cliente de manera automatizada, optimizando el tiempo del personal técnico para sus labores operativas. \\
\hline
\end{tabular}
\end{table}
\section{Objetivos}

\subsection{Objetivo General}

Diseñar y desarrollar un sistema de gestión y trazabilidad de servicios de soporte técnico con asistencia mediante inteligencia artificial para los prestadores de servicios de la ciudad de Tarija durante las gestiones 2026--2027, a partir del análisis de los procesos actuales y del levantamiento de requerimientos, con la finalidad de centralizar el registro de las órdenes de trabajo y facilitar su seguimiento y la comunicación con los clientes.

\subsection{Objetivos Específicos}

\begin{enumerate}
    \item Analizar el proceso actual de recepción de equipos, registro de servicios y comunicación con los clientes en prestadores de soporte técnico de la ciudad de Tarija, identificando las principales necesidades de información.
    \item Determinar los requerimientos funcionales y no funcionales del sistema, considerando el registro de datos de recepción, la gestión de estados del servicio y las necesidades de consulta de los usuarios.
    \item Modelar la estructura del sistema, los flujos de trabajo de las órdenes de servicio y la arquitectura del software orientada al desarrollo de la solución.
    \item Diseñar la lógica de interacción de un asistente conversacional basado en inteligencia artificial para apoyar la atención de consultas frecuentes y la entrega de información sobre el estado de los servicios.
    \item Desarrollar el prototipo funcional (MVP) de la solución propuesta y evaluar su factibilidad operativa mediante criterios y métricas de calidad de software según la norma ISO/IEC 25010:2023 para validar el sistema de trazabilidad de soporte técnico.
\end{enumerate}
\section{Justificación}

\subsection{Justificación técnica}

Desde la perspectiva técnica, el proyecto propone organizar la información de los servicios de soporte mediante un modelo estructurado de datos y la definición clara de estados de trabajo. Asimismo, contempla el diseño de un asistente conversacional basado en inteligencia artificial para apoyar la atención de consultas frecuentes sobre el estado de los equipos. Esta herramienta actuará como un canal complementario de información, sin sustituir la revisión técnica ni el diagnóstico especializado a cargo del personal calificado.

\subsection{Justificación social}

En el ámbito social, la propuesta busca facilitar el acceso de los clientes a la información sobre sus servicios de soporte técnico, reduciendo la incertidumbre durante el periodo de reparación. El diseño del sistema podría favorecer una comunicación más clara entre el usuario y el prestador del servicio, promoviendo acuerdos transparentes sobre los trabajos requeridos.

\subsection{Justificación económica}

Desde el punto de vista económico, el diseño propuesto busca respaldar la gestión organizativa de los prestadores de servicio. La centralización de la información y la automatización de respuestas a consultas frecuentes podrían contribuir a reducir el tiempo dedicado a la atención de llamadas o mensajes repetitivos, permitiendo un uso más eficiente del tiempo de trabajo.

\subsection{Justificación medioambiental}

En el aspecto ambiental, la digitalización de los registros de recepción y del seguimiento de las órdenes de trabajo podría reducir progresivamente la necesidad de utilizar comprobantes en papel. Adicionalmente, el adecuado mantenimiento de los equipos informáticos constituye una práctica que contribuye a prolongar su vida útil, lo que representa un beneficio indirecto potencial en favor del aprovechamiento de los recursos tecnológicos.

\section{Metodología de la investigación}

La metodología de la investigación orienta el proceso de recolección de datos y el análisis de información necesarios para caracterizar la situación actual del soporte técnico en la ciudad de Tarija y fundamentar la propuesta de diseño.

\section{Enfoque de la investigación}

El estudio adopta un enfoque mixto (cualitativo y cuantitativo). El componente cualitativo permitirá indagar sobre las experiencias, modalidades de trabajo y dificultades expresadas por los prestadores de servicio y los clientes mediante entrevistas semiestructuradas. El componente cuantitativo facilitará la recolección de datos descriptivos sobre la frecuencia de servicios, tipos de consultas más habituales y apreciaciones sobre los medios de comunicación utilizados mediante encuestas estructuradas. Ambos conjuntos de datos se articularán mediante una triangulación metodológica, orientada a traducir los hallazgos operativos y las necesidades de los usuarios en la definición formal de los requerimientos funcionales y no funcionales del sistema.
\section{Nivel de investigación}

La investigación es de nivel descriptivo y propositivo. Es descriptiva porque busca caracterizar cómo se realizan actualmente las actividades de recepción, seguimiento y entrega en los servicios de soporte técnico en el contexto local. Es propositiva porque, a partir del diagnóstico elaborado, formula el diseño y desarrollo de un sistema de gestión y trazabilidad como alternativa de solución informática, materializada en un prototipo funcional (MVP) evaluado según criterios de calidad de software.
\section{Tipo de estudio}

El estudio presenta un diseño no experimental y transversal. Es no experimental debido a que se observará la realidad de los prestadores de servicio sin manipular variables. Es transversal porque la recolección de información en el campo se llevará a cabo en un periodo determinado durante la gestión 2026.

\section{Métodos de investigación}

Se emplearán métodos teóricos como el análisis y la síntesis para revisar la documentación y fundamentar el diseño del sistema, así como el enfoque sistémico para comprender la interacción entre los componentes del software. Como métodos empíricos se utilizarán la observación, la entrevista y la encuesta para el levantamiento de información en el trabajo de campo.

\section{Población, tipo de muestreo y muestra}

La población de estudio estará conformada por prestadores de servicios de soporte técnico de la ciudad de Tarija y por usuarios que acuden a estos servicios. El muestreo será no probabilístico por conveniencia, seleccionando un grupo de talleres y clientes dispuestos a participar en las entrevistas y encuestas de diagnóstico.

\section{Técnicas e instrumentos de recolección de datos}

A continuación, se detalla la correspondencia entre los métodos empíricos, las técnicas, los instrumentos y la finalidad de la recolección de datos durante el diagnóstico de campo:

\begin{table}[H]
\centering
\caption{Correspondencia entre métodos, técnicas, instrumentos y finalidad}
\label{tab:tecnicas}
\small
\begin{tabular}{p{2.5cm} p{3cm} p{3cm} p{5cm}}
\toprule
\textbf{Método Empírico} & \textbf{Técnica de Recolección} & \textbf{Instrumento Aplicado} & \textbf{Finalidad de la Recolección} \\
\midrule
Método Empírico (Cualitativo) & Entrevista semiestructurada & Guía de preguntas abiertas en profundidad & Comprender los procesos internos de atención, registro de datos y principales dificultades de comunicación en taller. \\
\addlinespace
Método Empírico (Cuantitativo) & Encuesta estructurada & Cuestionario estructurado (físico y Google Forms) & Cuantificar patrones de consulta, medios preferidos de comunicación y percepción de incertidumbre en clientes. \\
\addlinespace
Método Empírico (Documental) & Análisis y revisión documental & Ficha de registro y matriz de cotejo de comprobantes & Identificar los datos mínimos requeridos para la trazabilidad legal e histórica de las órdenes de servicio. \\
\bottomrule
\end{tabular}
\end{table}

\textbf{Procedimiento de Análisis de Datos:} Los datos cuantitativos recolectados mediante encuestas estructuradas se procesarán con estadística descriptiva (distribución de frecuencias y porcentajes) mediante hojas de cálculo (Excel), generando representaciones gráficas para caracterizar los patrones de consulta. Los datos cualitativos provenientes de las entrevistas semiestructuradas se analizarán mediante análisis temático de contenido, identificando categorías recurrentes sobre los problemas de registro y comunicación. La información resultante se integrará mediante triangulación metodológica para definir formalmente la especificación de requerimientos de software según la norma ISO/IEC/IEEE 29148:2018.

\textbf{Consideraciones Éticas:} La investigación se rige por principios de ética científica, garantizando la participación estrictamente voluntaria y el consentimiento informado previo de técnicos y clientes. Se aplica una política de minimización de datos y anonimización estricta, protegiendo la identidad de los participantes y el secreto comercial de los talleres en cumplimiento del Artículo 21 (inciso 2) y Artículo 130 de la Constitución Política del Estado de Bolivia y la Ley N.º 164.
\section{Medios}

Para la realización del proyecto se utilizarán recursos de software para el procesamiento de datos, herramientas de modelado para la estructuración del sistema y formularios digitales para la aplicación de encuestas, complementados con equipos informáticos de oficina.

\section{Fuentes de información}

Se recurrirá a fuentes primarias mediante la información recolectada directamente de los participantes durante el diagnóstico de campo. Como fuentes secundarias se consultará literatura académica, textos especializados en ingeniería de software y guías metodológicas institucionales.


\chapter{Marco teórico}
\section{Marco histórico}

Pendiente de desarrollo.

\section{Marco Contextual}

El marco contextual delimita el entorno geográfico, económico, tecnológico e institucional en el cual se inserta la problemática investigada, permitiendo caracterizar tanto las tendencias observadas a nivel global como las particularidades operativas del escenario local en el que se proyecta la propuesta (Hernández Sampieri, 2014).

\subsection{Contexto internacional del soporte técnico y la trazabilidad de servicios TI}

En el ámbito internacional de las economías industrializadas, la prestación de servicios de soporte técnico se encuentra altamente automatizada mediante plataformas corporativas de Enterprise Resource Planning (ERP) y módulos especializados de gestión de servicios, tales como SAP Service Cloud (SAP SE, 2023). Asimismo, en el segmento de pequeñas y medianas empresas se ha masificado el uso de soluciones modulares e integradas de código abierto o modelo freemium/SaaS como Odoo, la cual incluye módulos de mesa de ayuda (Helpdesk), inventarios y reparación (Odoo S.A., 2024). Estas herramientas demuestran que la trazabilidad integral incrementa la satisfacción del usuario y reduce los costos operativos de búsqueda documental. No obstante, el análisis de sus funcionalidades evidencia que plataformas como SAP o los módulos avanzados de Odoo requieren costos de licenciamiento en divisas, procesos complejos de configuración y administración continua por personal de TI dedicado, características que las vuelven poco idóneas o costosas para talleres independientes de reparación de computadoras en entornos emergentes.

\subsection{Contexto latinoamericano y brechas operativas en micro y pequeñas empresas}

En el contexto latinoamericano, si bien las empresas medianas y grandes han adoptado plataformas corporativas de soporte e infraestructura administrada a través de proveedores regionales como Sonda o Noventiq (Sonda S.A., 2024; Noventiq, 2024), existe una marcada brecha tecnológica en las micro y pequeñas empresas (MIPYMES) de mantenimiento informático. Estudios como el de Mosquera Perea et al. (2025) en Colombia evidencian que los talleres independientes de reparación sufren de desorganización en el registro de recepción y falta de mecanismos para demostrar la transparencia en las cotizaciones. Asimismo, soluciones de código abierto populares en la región como GLPI (GLPI Project, 2024), si bien destacan en la gestión de inventarios y tickets de TI corporativos internos, carecen de canales conversacionales con inteligencia artificial orientados al cliente final y de flujos ágiles de cotización y evidencias fotográficas para talleres comerciales independientes.

\subsection{Contexto nacional (Bolivia) y madurez del ecosistema digital}

En Bolivia, el ecosistema de base digital formal presenta una marcada concentración geográfica en las ciudades del eje troncal: Santa Cruz (40,9\%), La Paz (33,8\%) y Cochabamba (20,1\%), acumulando el 94,8\% del total nacional de emprendimientos digitales registradas en el reporte del Mapeo del Ecosistema de Tecnología Digital en Bolivia 2026 (Fundación Emprender Futuro et al., 2026). A nivel nacional existen instituciones promotoras de la trazabilidad empresarial como GS1 Bolivia (GS1 Bolivia, 2026), la cual brinda estándares de identificación y código de barras. Asimismo, en la gestión de servicios informáticos destacan soluciones Open Source como GLPI (GLPI Project, 2024) y software de gestión especializado como ExpertoSis, Freshservice o Jira Service Management. No obstante, estas plataformas están orientadas a inventarios institucionales o mesas de ayuda corporativas y no a la dinámica comercial directa entre un taller técnico independiente y sus clientes particulares.

\subsection{Contexto local (Ciudad de Tarija) y dinámicas operativas de los talleres independientes}

En el radio urbano de la ciudad de Tarija, la presente investigación se orienta al estudio de los prestadores de servicios de soporte técnico y mantenimiento de computadoras, abarcando tanto talleres comerciales establecidos como iniciativas y técnicos independientes. Como parte del diagnóstico de campo planificado, se indagará de manera directa sobre las modalidades operativas utilizadas para registrar la recepción de equipos, documentar los diagnósticos, comunicar los presupuestos y dar seguimiento a las reparaciones.

En este marco, se examinará empíricamente el grado de adopción de registros físicos en papel, herramientas digitales locales o aplicaciones de mensajería informal, evitando asumir anticipadamente la prevalencia de una modalidad específica. Asimismo, la ausencia de un directorio gremial o registro comercial centralizado se considera una condición del entorno que justifica la necesidad de realizar un diagnóstico directo. Los hallazgos del trabajo de campo permitirán fundamentar la definición de requerimientos del sistema propuesto en las dinámicas reales identificadas en el contexto tarijeño.

A continuación, se presenta un cuadro comparativo entre las principales soluciones tecnológicas de gestión de servicios y la propuesta de diseño formulada para la ciudad de Tarija:

\begin{table}[H]
\centering
\footnotesize
\setlength{\tabcolsep}{3pt}
\renewcommand{\arraystretch}{1.15}
\begin{tabular}{|p{2.3cm}|p{2cm}|p{2cm}|p{2.9cm}|p{2.5cm}|}
\hline
\textbf{Solución / Plataforma} & \textbf{Licencia / Modelo} & \textbf{Dominio Principal} & \textbf{Trazabilidad e IA} & \textbf{Adecuación a Talleres Tarija} \\
\hline
Odoo (Helpdesk/Repair) & Freemium / SaaS Modulares & ERP General para PYMES & Trazabilidad de tickets; IA mediante módulos pagados. & Parcial: Exige configuración compleja de múltiples módulos. \\
\hline
SAP Service Cloud & De Paga (Corporativa) & Gran Empresa y Corporaciones & Trazabilidad avanzada; IA predictiva empresarial. & Baja: Costos de licenciamiento inaccesibles para talleres pyme. \\
\hline
GLPI (Open Source) & Código Abierto (GPLv3) & Gestión TI e Inventarios de Red & Trazabilidad por tickets; IA vía plugins externos. & Media: Enfoque en departamentos TI internos, no en cliente final. \\
\hline
GS1 Bolivia & Estándares Internacionales & Trazabilidad Logística y Código de Barras & Estándares de identificación y códigos GS1; sin enfoque en tickets o IA. & Baja: Orientado a cadenas de suministro e industria general, no a talleres informáticos. \\
\hline
ExpertoSis (Sistemas de Gestión TI) & Comercial (Propietario) & Gestión de Incidencias y Servicios TI & Gestión de tickets y módulo de trazabilidad (según especificaciones comerciales); sin IA conversacional. & Media: Enfocado en soporte empresarial genérico, sin asistentes conversacionales desacoplados. \\
\hline
Sistema Propuesto (Tarija) & Licencia Propia (Diseño a Medida) & Soporte Técnico de Computadoras & Diseño de trazabilidad por estados secuenciales y evidencias fotográficas; propuesta de asistente de IA desacoplado para consultas de estado e información. & Alta: Diseñado específicamente para las dinámicas de los talleres independientes de Tarija. \\
\hline
\end{tabular}
\caption{Cuadro comparativo de soluciones informáticas de gestión y trazabilidad de soporte técnico.}
\label{tab:comparativo-soluciones}
\end{table}
\section{Marco conceptual}

Pendiente de desarrollo.

\section{Marco Normativo}

El marco normativo identifica y analiza las disposiciones legales vigentes en el Estado Plurinacional de Bolivia y los estándares técnicos internacionales aplicables al diseño del sistema. Dado que la plataforma gestionará datos de identificación personal de usuarios, registros de activos informáticos y acuerdos comerciales de servicio, resulta imperativo que la solución tecnológica se alinee con los principios constitucionales de privacidad y los estándares de calidad y seguridad de la industria del software.

\subsection{Marco Jurídico Nacional (Bolivia)}

El ordenamiento jurídico boliviano establece las garantías fundamentales para la protección de la información ciudadana y la formalización de las relaciones contractuales y comerciales en entornos digitales:

\begin{itemize}
    \item \textbf{Constitución Política del Estado Plurinacional de Bolivia:} El Artículo 21 (inciso 2) consagra el derecho fundamental de toda persona a la privacidad, intimidad, honra, propia imagen y dignidad. Asimismo, el Artículo 130 establece la Acción de Protección de Privacidad, otorgando a toda persona la facultad de conocer, actualizar o rectificar los datos personales registrados en bancos de datos públicos o privados. El diseño del software incorpora este mandato al restringir la visibilidad de los datos personales del cliente exclusivamente al personal autorizado del taller.
    \item \textbf{Ley N.º 164 (Ley General de Telecomunicaciones y TIC):} Promulgada el 8 de agosto de 2011, esta ley establece el marco regulatorio general para las tecnologías de la información y comunicación. Sus Artículos 6 y 54 reconocen la inviolabilidad y el secreto de las comunicaciones electrónicas, así como la plena validez jurídica y probatoria de los mensajes de datos, documentos electrónicos y firmas digitales. La propuesta cumple esta norma al formalizar la emisión del comprobante digital de servicio y el registro auditable de autorizaciones del cliente.
    \item \textbf{Ley N.º 453 (Ley General de los Derechos de las Usuarias y los Usuarios y de las Consumidoras y los Consumidores):} Promulgada el 4 de diciembre de 2013, esta norma regula las relaciones entre proveedores de servicios y consumidores. En su Artículo 11 (incisos 1, 2 y 3) consagra el derecho de las usuarias y los usuarios a recibir información fidedigna, veraz, oportuna, clara y completa sobre las características, precios y condiciones de los servicios contratados, así como a un trato equitativo y sin discriminación. El sistema propuesto respalda directamente este derecho al ofrecer a los clientes acceso transparente y continuo al estado de su equipo y al detalle pormenorizado de las intervenciones técnicas realizadas.
\end{itemize}

\subsection{Estándares Técnicos Internacionales de Calidad y Seguridad}

En el ámbito de la ingeniería de software, la aplicación de normas internacionales asegura que el artefacto diseñado satisfaga criterios rigurosos de calidad estructural y operativa:

\begin{itemize}
    \item \textbf{Estándar ISO/IEC 25010:2023 (Modelos de Calidad del Producto Software):} Pertenece a la familia SQuaRE (Software Product Quality Requirements and Evaluation) y en su edición de 2023 define un modelo para evaluar la calidad del producto software a través de nueve características principales: adecuación funcional (functional suitability), eficiencia de desempeño (performance efficiency), compatibilidad (compatibility), capacidad de interacción (interaction capability), fiabilidad (reliability), seguridad (security), mantenibilidad (maintainability), flexibilidad (flexibility) y protección frente a riesgos (safety). Para la especificación de requerimientos no funcionales del sistema propuesto, se priorizan la adecuación funcional, la capacidad de interacción, la fiabilidad y la seguridad informacional.
    \item \textbf{Estándares ISO/IEC 27001 e ISO/IEC 27002 (Seguridad de la Información):} La norma ISO/IEC 27001 especifica los requisitos para establecer, implementar, mantener y mejorar un Sistema de Gestión de la Seguridad de la Información (SGSI), mientras que la norma ISO/IEC 27002 actúa como una guía de buenas prácticas que describe controles de seguridad aplicables a la gestión de accesos, protección de contraseñas mediante funciones criptográficas unidireccionales (hashing con algoritmos seguros como Bcrypt), generación de copias de respaldo de la base de datos y políticas de auditoría interna de transacciones.
\end{itemize}

\subsubsection{Matriz Operativa de Evaluación de Calidad del MVP (ISO/IEC 25010:2023)}

Para operacionalizar la evaluación del prototipo funcional (MVP) de acuerdo con el Objetivo Específico 5, la propuesta establece una matriz de métricas e indicadores de calidad fundamentada en el estándar ISO/IEC 25010:2023, la cual se aplicará en el entorno de pruebas durante Taller de Grado II:

\begin{table}[H]
\centering
\caption{Matriz operativa de evaluación de calidad del MVP según ISO/IEC 25010:2023}
\label{tab:matriz-iso25010}
\small
\begin{tabular}{p{2.2cm} p{2.5cm} p{3cm} p{3cm} p{2.3cm}}
\toprule
\textbf{Dimensión de Calidad} & \textbf{Indicador Técnico} & \textbf{Método de Verificación / Cálculo} & \textbf{Instrumento / Entorno} & \textbf{Criterio de Aceptación} \\
\midrule
Adecuación Funcional & Cobertura de Casos de Prueba & (Casos de prueba aprobados / Ejecutados) $\times$ 100 & Script de pruebas funcionales ($\sim$20 flujos núcleo: Recepción, Diagnóstico, Autorización, Reparación, Pruebas, Entrega, Asistente IA, RBAC) & $\geq$ 95\% casos aprobados \\
\addlinespace
Capacidad de Interacción & Puntuación de Usabilidad Percibida & Cuestionario System Usability Scale (SUS) & Escala SUS a muestra piloto (20 participantes: 5 técnicos, 15 clientes) & Puntaje SUS $\geq$ 75/100 \\
\addlinespace
Fiabilidad (Confiabilidad) & Tasa de Fallos Transaccionales & (Transacciones fallidas / Probadas) $\times$ 100 & Logs de aplicación y pruebas de estrés (50 transacciones simuladas) & 0\% fallos en transiciones \\
\addlinespace
Seguridad de la Información & Efectividad de Controles RBAC & (Peticiones no autorizadas bloqueadas / Intentos) $\times$ 100 & Pruebas de penetración y auditoría (30 escenarios: RBAC, UUIDv4, prompt injection) & 100\% intentos bloqueados \\
\addlinespace
Eficiencia de Desempeño (Backend) & Tiempo de Respuesta HTTP Backend & Tiempo promedio de respuesta HTTP en API REST interna & Herramienta ApacheBench / JMeter sobre servidor Laravel & Latencia HTTP $\leq$ 2,0 s \\
\addlinespace
Eficiencia de Desempeño (Asistente IA) & Tiempo de Respuesta API del LLM & Medición del tiempo total de respuesta de la API externa & Logs de peticiones REST a Google Gemini API & Latencia API IA $\leq$ 4,0 s \\
\bottomrule
\end{tabular}
\end{table}
\section{Marco Referencial (Pila Tecnológica del Sistema)}

El marco referencial o tecnológico sistematiza las herramientas, plataformas, lenguajes de programación, marcos de trabajo (frameworks) y servicios de infraestructura informática seleccionados para sustentar el diseño de la arquitectura del sistema de gestión y trazabilidad. A diferencia del marco conceptual ---orientado a los principios y modelos teóricos abstractos---, este apartado especifica la viabilidad técnica y las características operativas de la pila tecnológica (stack) propuesta, orientada a respaldar una articulación eficiente entre la persistencia relacional, la lógica de negocio y la integración con modelos de inteligencia artificial.

\subsection{Entorno de Desarrollo y Lenguajes de Programación}

\begin{itemize}
    \item \textbf{PHP \cite{php2024}:} Lenguaje de programación de código abierto del lado del servidor ampliamente utilizado en aplicaciones web corporativas y sistemas transaccionales. Su compatibilidad nativa con sistemas gestores de bases de datos relacionales, su soporte robusto para Programación Orientada a Objetos (POO) y su amplia adopción en arquitecturas de soporte técnico (como se evidencia en los antecedentes de Mosquera Perea et al., \cite{mosquera2025}, y Velasco Castro, \cite{velasco2023}) lo constituyen como el núcleo para la implementación de la lógica del servidor.
    \item \textbf{JavaScript / TypeScript \cite{ecma2023}:} Lenguajes de programación estándar para la ejecución del lado del cliente en navegadores web. Permiten gestionar la asincronía en las solicitudes de estado, dinamizar la carga de interfaces reactivas y procesar la interacción conversacional del usuario en tiempo real sin recargar la página completa.
    \item \textbf{HTML5 y CSS3 \cite{w3c2023}:} Estándares fundamentales para la estructuración semántica del contenido digital y la definición de hojas de estilo en cascada. Garantizan la correcta presentación visual, la accesibilidad de la información de las órdenes de servicio y la compatibilidad multiplataforma en dispositivos de escritorio y móviles.
\end{itemize}

\subsection{Marcos de Trabajo (Frameworks) y Librerías de Interfaz}

\begin{itemize}
    \item \textbf{Laravel \cite{otwell2024}:} Framework de desarrollo web estructurado bajo el patrón arquitectónico Modelo-Vista-Controlador (MVC). Laravel provee mecanismos nativos esenciales para el proyecto: enrutamiento estructurado, motor de mapeo objeto-relacional (ORM Eloquent) para la gestión transparente de consultas a la base de datos, sistema de autenticación y autorización robusto, middleware para la protección contra vulnerabilidades comunes (como inyección SQL y CSRF), y facilidades para la construcción de APIs RESTful que interconecten el backend con servicios externos.
    \item \textbf{Bootstrap 5 \cite{otto2023}:} Framework de diseño frontend enfocado en la creación de interfaces responsivas bajo el paradigma mobile-first. Permite estructurar tableros de control (dashboards) intuitivos y adaptables a pantallas de computadoras, tabletas y teléfonos móviles, facilitando que tanto el personal técnico en taller como los clientes desde sus dispositivos personales puedan interactuar con la plataforma sin barreras visuales.
\end{itemize}

\subsection{Motor de Base de Datos y Persistencia Relacional}

\begin{itemize}
    \item \textbf{MySQL 8.0 \cite{oracle2024}:} Sistema de Gestión de Bases de Datos Relacionales (RDBMS) de alto rendimiento, confiabilidad y soporte para transacciones ACID (Atomicidad, Consistencia, Aislamiento y Durabilidad). Se selecciona para el diseño del esquema entidad-relación del sistema, orientado a respaldar la integridad referencial entre clientes, dispositivos, diagnósticos, órdenes de trabajo, técnicos asignados y marcas de tiempo (timestamps). Asimismo, contempla implementar restricciones de borrado lógico (soft deletes) a nivel de ORM y un esquema de auditoría mediante tablas de registro (audit\_logs) y permisos restrictivos de usuario DBMS (revocación de permisos DELETE), orientados a fortalecer la integridad del sistema y preservar un historial auditable reduciendo el riesgo de alteraciones no autorizadas.
    \item \textbf{MySQL Workbench:} Herramienta visual de modelado, administración y diseño de bases de datos que facilita la creación de esquemas físicos, diagramas entidad-relación y ejecución controlada de scripts de inicialización y migración.
\end{itemize}

\subsection{Herramientas de Modelado de Software y Documentación}

\begin{itemize}
    \item \textbf{Enterprise Architect / StarUML \cite{omg2017}:} Herramientas especializadas de modelado basadas en el estándar UML 2.5. Se emplearán para elaborar los cinco diagramas formales de ingeniería de software contemplados en la arquitectura del sistema: Diagramas de Casos de Uso (definición de roles y límites del sistema), Diagramas de Actividades (flujos de trabajo del ciclo de reparación), Diagramas de Clases (modelo de dominio y persistencia), Diagramas de Secuencia (intercambio temporal de mensajes) y Diagramas de Despliegue (topología física de hardware, servidor web Laravel, RDBMS MySQL y API de IA).
\end{itemize}

\subsection{Pila Tecnológica de Inteligencia Artificial Conversacional}

\begin{itemize}
    \item \textbf{Modelo de Lenguaje de Gran Tamaño en la Nube (Google Gemini API REST):} Se establece como alternativa tecnológica de referencia la plataforma Google Gemini \cite{googleai2026}, proyectada para su integración mediante API REST con respuestas estructuradas en formato JSON, cuya selección y versión definitiva de modelo se evaluará en función de su disponibilidad, estabilidad de servicio, latencia de procesamiento y rendimiento en español.
    \item \textbf{Función de Interpretación y Redacción del LLM:} En el diseño propuesto, el modelo de lenguaje actuará estrictamente como motor de comprensión semántica e interfaz conversacional en español. No poseerá permisos de ejecución de código ni acceso directo a las tablas del sistema.
    \item \textbf{Mecanismo de Desacoplamiento y Control de Acceso (API REST HTTPS Intermediaria):} Capa de software intermedia proyectada en el backend de Laravel para validar el token verificado de la orden de servicio, realizar lecturas seguras de solo lectura en MySQL y proveer los datos necesarios al LLM dentro de un entorno controlado mediante ingeniería de instrucciones (system prompts).
\end{itemize}

\nocite{*}
\renewcommand{\bibname}{REFERENCIAS BIBLIOGRÁFICAS}
\bibliographystyle{apalike}
\bibliography{bibliografia}

\renewcommand{\tablename}{Tabla}
\renewcommand{\listtablename}{Índice de tablas}

\end{document}